\documentclass[12pt,a4paper]{article}

% ── Packages ────────────────────────────────────────────────────────────────
\usepackage{amsmath, amssymb, bm}
\usepackage{geometry}
\usepackage{booktabs}
\usepackage{array}
\usepackage{xcolor}
\usepackage{hyperref}
\usepackage{parskip}
\usepackage{titlesec}

\geometry{margin=2.5cm}
\hypersetup{colorlinks=true, linkcolor=blue, urlcolor=blue}

% ── Convenience macros ──────────────────────────────────────────────────────
\newcommand{\cph}{\cos\phi}
\newcommand{\sph}{\sin\phi}
\newcommand{\cth}{\cos\theta}
\newcommand{\sth}{\sin\theta}
\newcommand{\cps}{\cos\psi}
\newcommand{\sps}{\sin\psi}
\newcommand{\tph}{\tan\phi}
\newcommand{\tth}{\tan\theta}
\newcommand{\q}{\mathbf{q}}
\newcommand{\bu}{\mathbf{u}}
\newcommand{\bF}{\mathbf{F}}
\newcommand{\bP}{\mathbf{P}}
\newcommand{\bK}{\mathbf{K}}
\newcommand{\bH}{\mathbf{H}}
\newcommand{\bQ}{\mathbf{Q}}
\newcommand{\bR}{\mathbf{R}}
\newcommand{\bI}{\mathbf{I}}
\newcommand{\eone}{\mathbf{e}_1}

% ── Title ────────────────────────────────────────────────────────────────────
\title{%
  \textbf{Mathematical Foundations}\\[0.4em]
  \large Autonomous Frontier Exploration with Semantic Hazard Mapping\\[0.2em]
  \normalsize RAS 598: Mobile Robotics $\cdot$ Spring 2026 $\cdot$ ASU
}
\author{Rohit Mane \texttt{(rmane2)}}
\date{}

% ══════════════════════════════════════════════════════════════════════════════
\begin{document}
\maketitle
\tableofcontents
\newpage

% ══════════════════════════════════════════════════════════════════════════════
\section{2D Differential Drive Kinematics}
% ══════════════════════════════════════════════════════════════════════════════

\subsection{State Vector and Control Inputs}

The TurtleBot4 Lite uses an iRobot Create3 differential drive base.
The robot's pose is described by three quantities.
The control inputs are the angular velocities of the left and right wheels:

\begin{equation}
  \q(t) = \begin{bmatrix} x(t) \\ y(t) \\ \psi(t) \end{bmatrix},
  \qquad
  \bu(t) = \begin{bmatrix} \omega_L(t) \\ \omega_R(t) \end{bmatrix}
  \tag{1}
\end{equation}

\begin{center}
\begin{tabular}{>{\(}l<{\)} l l}
  \toprule
  \text{Symbol} & Meaning & Value (TurtleBot4 Lite) \\
  \midrule
  x,\, y        & Position of robot centre in world frame & --- \\
  \psi          & Heading angle from positive $x$-axis [rad] & --- \\
  \omega_L,\,\omega_R & Left/right wheel angular velocities [rad\,s$^{-1}$] & max 8.8 rad\,s$^{-1}$ \\
  \bottomrule
\end{tabular}
\end{center}

% ──────────────────────────────────────────────────────────────────────────────
\subsection{Wheel Velocities to Linear and Angular Velocity}

Each wheel of radius $r$ contributes a tangential speed. The net linear speed $v$
and yaw rate $\dot\psi$ are their symmetric and antisymmetric combinations,
normalised by the wheelbase $L$:

\begin{equation}
  v(t) = \frac{r}{2}\bigl(\omega_R(t) + \omega_L(t)\bigr)
  \tag{2}
\end{equation}

\begin{equation}
  \dot{\psi}(t) = \frac{r}{L}\bigl(\omega_R(t) - \omega_L(t)\bigr)
  \tag{3}
\end{equation}

\begin{center}
\begin{tabular}{>{\(}l<{\)} l l}
  \toprule
  \text{Symbol} & Meaning & TurtleBot4 Lite value \\
  \midrule
  r & Wheel radius [m] & 0.0352\,m \\
  L & Track width (wheel-to-wheel) [m] & 0.233\,m \\
  v & Forward linear speed [m\,s$^{-1}$] & max 0.31\,m\,s$^{-1}$ \\
  \dot\psi & Yaw rate [rad\,s$^{-1}$] & max 1.90\,rad\,s$^{-1}$ \\
  \bottomrule
\end{tabular}
\end{center}

% ──────────────────────────────────────────────────────────────────────────────
\subsection{Continuous-Time Motion Model (ODE)}

Projecting the forward speed onto the world frame gives the full nonlinear
kinematic model:

\begin{equation}
  \dot{\q}(t)
  = \begin{bmatrix} \dot{x}(t) \\ \dot{y}(t) \\ \dot{\psi}(t) \end{bmatrix}
  = \begin{bmatrix} v(t)\cos\psi(t) \\ v(t)\sin\psi(t) \\ \dot{\psi}(t) \end{bmatrix}
  = f\bigl(\q(t),\,\bu(t)\bigr)
  \tag{4}
\end{equation}

% ──────────────────────────────────────────────────────────────────────────────
\subsection{Discrete-Time State Update (Euler Integration)}

For a fixed timestep $\Delta t$ using forward Euler:

\begin{equation}
  \q_{k+1} = \q_k + \Delta t \cdot f(\q_k,\,\bu_k)
  \tag{5}
\end{equation}

Expanding the three state components explicitly:

\begin{equation}
  \begin{aligned}
    x_{k+1}   &= x_k   + v_k \cos\psi_k \cdot \Delta t \\
    y_{k+1}   &= y_k   + v_k \sin\psi_k \cdot \Delta t \\
    \psi_{k+1} &= \psi_k + \dot{\psi}_k  \cdot \Delta t
  \end{aligned}
  \tag{6}
\end{equation}

where $v_k$ and $\dot\psi_k$ are computed from
$(\omega_{L,k},\,\omega_{R,k})$ via equations~(2)--(3).

% ──────────────────────────────────────────────────────────────────────────────
\subsection{Exact Arc-Based Integration (Curved Motion)}

When $\dot\psi_k \neq 0$, integrating exactly over the arc avoids
Euler accumulation error. Let $\kappa = \dot\psi_k\,\Delta t$:

\begin{equation}
  \begin{aligned}
    x_{k+1}   &= x_k + \frac{v_k}{\dot{\psi}_k}
                  \Bigl(\sin(\psi_k+\kappa) - \sin\psi_k\Bigr) \\
    y_{k+1}   &= y_k - \frac{v_k}{\dot{\psi}_k}
                  \Bigl(\cos(\psi_k+\kappa) - \cos\psi_k\Bigr) \\
    \psi_{k+1} &= \psi_k + \kappa
  \end{aligned}
  \tag{7}
\end{equation}

When $\dot\psi_k \approx 0$ (straight-line motion), equation~(6) is
numerically preferred; equation~(7) degenerates gracefully via
L'H\^{o}pital's rule to the same result.

% ──────────────────────────────────────────────────────────────────────────────
\subsection{Full Mapping: Control Inputs $\to$ State Update}

Combining equations~(2)--(3) and~(6), the explicit mapping from
wheel velocities to state update is:

\begin{equation}
  \begin{bmatrix} x_{k+1} \\ y_{k+1} \\ \psi_{k+1} \end{bmatrix}
  =
  \begin{bmatrix} x_k \\ y_k \\ \psi_k \end{bmatrix}
  +
  \frac{r\,\Delta t}{2}
  \begin{bmatrix}
    (\omega_{R,k}+\omega_{L,k})\cos\psi_k \\
    (\omega_{R,k}+\omega_{L,k})\sin\psi_k \\
    \dfrac{2}{L}(\omega_{R,k}-\omega_{L,k})
  \end{bmatrix}
  \tag{8}
\end{equation}

% ══════════════════════════════════════════════════════════════════════════════
\newpage
\section{Frontier Scoring Formula}
% ══════════════════════════════════════════════════════════════════════════════

\subsection{The Formula}

Among all candidate frontier regions $R$ on the boundary between
known-free and unknown space, the robot selects the one with the highest score:

\begin{equation}
  \text{score}(R)
  = \frac{\alpha \cdot \text{size}(R) \;+\; \beta \cdot \text{unknown}(R)}
         {\text{dist}(R) + \varepsilon}
  \tag{9}
\end{equation}

\begin{center}
\begin{tabular}{>{\(}l<{\)} l l}
  \toprule
  \text{Symbol} & Meaning & Value in code \\
  \midrule
  \text{size}(R)    & Number of frontier cells in region $R$ & cluster size \\
  \text{unknown}(R) & Unknown cells within $r{=}0.5$\,m of centroid & count at radius \\
  \text{dist}(R)    & Euclidean distance from robot to centroid [m] & metres \\
  \alpha            & Size weight & \texttt{1.0} \\
  \beta             & Information gain weight & \texttt{2.0} \\
  \varepsilon       & Division-by-zero guard & \texttt{1e-6} \\
  \bottomrule
\end{tabular}
\end{center}

\noindent\textbf{Note on parameter values:}
The implemented values are $\alpha = 1.0$, $\beta = 2.0$, $\varepsilon = 10^{-6}$.
These weight information gain twice as heavily as frontier size
($\alpha/\beta = 0.5$), driving aggressive map coverage.
Only the ratio $\alpha/\beta$ affects frontier ranking ---
scaling both by a constant leaves all scores unchanged.

% ──────────────────────────────────────────────────────────────────────────────
\subsection{Parameter Rationale}

\paragraph{$\alpha$ --- size weight (1.0).}
Favours large open frontiers over isolated pockets, preventing the robot
from wasting navigation budget on tiny dead-end corridors.
Typical range: $\alpha \in [0.3,\,1.0]$.

\paragraph{$\beta$ --- unknown weight (2.0).}
Directly rewards information gain. Higher $\beta$ drives aggressive exploration;
lower $\beta$ produces cautious, thorough coverage.
Typical range: $\beta \in [1.0,\,3.0]$.

\paragraph{$\varepsilon$ --- stability guard (\texttt{1e-6}).}
Must satisfy $\varepsilon \ll \min_R \text{dist}(R)$.
At $10^{-6}$\,m it stabilises the denominator without biasing scores for
any realistic robot position.

\paragraph{The $\alpha/\beta$ ratio.}
\begin{itemize}
  \item $\alpha/\beta < 1$: robot prioritises information density over area.
  \item $\alpha/\beta > 1$: robot prefers large open spaces.
\end{itemize}
Our ratio $\alpha/\beta = 0.5$ prioritises information density.

% ──────────────────────────────────────────────────────────────────────────────
\subsection{Failure Modes and Mitigations}

\begin{description}
  \item[Greedy cycling.]
        Two frontiers with equal scores cause oscillation.
        Mitigation: hysteresis bonus for currently-targeted frontier.

  \item[Distance dominance late in mapping.]
        When all remaining unknowns are far away, the denominator overwhelms
        the numerator and scores compress.
        Mitigation: normalise $\text{dist}(R)$ by the map diagonal,
        or use a log penalty $\ln(\text{dist}+1)$.

  \item[$\alpha/\beta$ wrong for environment.]
        Open warehouses favour high $\alpha$; tight corridors favour high $\beta$.
        Mitigation: tune offline on a representative map.

  \item[$\varepsilon$ too large.]
        If $\varepsilon \gg \text{dist}(R)$ for near frontiers,
        the distance penalty vanishes.
        Keep $\varepsilon \ll$ minimum expected distance.
\end{description}

% ══════════════════════════════════════════════════════════════════════════════
\newpage
\section{SLAM: Occupancy Grid via Bayes Rule}
% ══════════════════════════════════════════════════════════════════════════════

SLAM Toolbox maintains a probabilistic occupancy grid where each cell $m_i$
stores the log-odds of being occupied given all sensor observations $z_{1:t}$.

\subsection{Log-Odds Representation}

The log-odds of a probability $p$ is defined as:

\begin{equation}
  l(p) = \log\left(\frac{p}{1-p}\right)
  \tag{10}
\end{equation}

This maps probabilities $[0,1]$ to $(-\infty,+\infty)$, making Bayesian updates
an additive operation that is numerically stable and computationally cheap.

\subsection{Recursive Bayesian Update}

\begin{equation}
  l(m_i \mid z_{1:t})
  = l(m_i \mid z_{1:t-1}) + l(m_i \mid z_t) - l_0
  \tag{11}
\end{equation}

\begin{center}
\begin{tabular}{>{\(}l<{\)} l l}
  \toprule
  \text{Symbol} & Meaning & Default value \\
  \midrule
  l(m_i \mid z_t) & Sensor model log-odds (LiDAR hit) & $+0.85$ \\
  l(m_i \mid z_t) & Sensor model log-odds (LiDAR miss) & $-0.4$ \\
  l_0             & Prior log-odds (uniform 50\% prior) & $0$ \\
  \bottomrule
\end{tabular}
\end{center}

To recover probability from log-odds:
$p = 1 - \dfrac{1}{1 + e^l}$.
The \texttt{OccupancyGrid} message encodes probabilities as integers
0--100 (0\,=\,free, 100\,=\,occupied, $-1$\,=\,unknown).

% ══════════════════════════════════════════════════════════════════════════════
\newpage
\section{Extended Kalman Filter: Prediction and Update}
% ══════════════════════════════════════════════════════════════════════════════

The \texttt{robot\_localization} EKF node fuses wheel odometry (\texttt{/odom})
and IMU data to produce \texttt{/odometry/filtered}, consumed by the frontier
explorer and behavior coordinator nodes.

\subsection{Prediction Step}

\begin{equation}
  \hat{\q}_{k+1|k}
  = \q_k + \Delta t
  \begin{bmatrix}
    v\cos\theta_k \\
    v\sin\theta_k \\
    \dfrac{r}{L}(\omega_R - \omega_L)
  \end{bmatrix}
  \tag{12}
\end{equation}

\begin{equation}
  \bP_{k+1|k} = \bF_k \bP_k \bF_k^\top + \bQ_k
  \tag{13}
\end{equation}

The Jacobian of the motion model (first-order Taylor linearisation):

\begin{equation}
  \bF_k = \bI + \Delta t
  \begin{bmatrix}
    0 & 0 & -v\sin\theta_k \\
    0 & 0 &  v\cos\theta_k \\
    0 & 0 & 0
  \end{bmatrix}
  \tag{14}
\end{equation}

\begin{center}
\begin{tabular}{>{\(}l<{\)} l}
  \toprule
  \text{Symbol} & Meaning \\
  \midrule
  \bP_k   & $3{\times}3$ state covariance matrix at step $k$ \\
  \bQ_k   & Process noise covariance (wheel slip, unmodelled dynamics) \\
  \bF_k   & Jacobian of $f$ with respect to state $\q$ \\
  \bottomrule
\end{tabular}
\end{center}

\subsection{Update Step}

When a sensor measurement $z_k$ arrives (IMU yaw rate, wheel encoder):

\begin{equation}
  \bK_k = \bP_{k|k-1}\bH_k^\top
  \left(\bH_k \bP_{k|k-1} \bH_k^\top + \bR_k \right)^{-1}
  \tag{15}
\end{equation}

\begin{equation}
  \hat{\q}_{k|k}
  = \hat{\q}_{k|k-1}
  + \bK_k \left(z_k - h(\hat{\q}_{k|k-1})\right)
  \tag{16}
\end{equation}

\begin{equation}
  \bP_{k|k} = (\bI - \bK_k \bH_k)\bP_{k|k-1}
  \tag{17}
\end{equation}

\begin{center}
\begin{tabular}{>{\(}l<{\)} l}
  \toprule
  \text{Symbol} & Meaning \\
  \midrule
  \bH_k & Jacobian of measurement model $h$ w.r.t.\ state \\
  \bR_k & Measurement noise covariance \\
  \bK_k & Kalman gain --- weights prediction vs.\ measurement trust \\
  z_k - h(\hat{\q}) & Innovation: difference between actual and predicted measurement \\
  \bottomrule
\end{tabular}
\end{center}

% ══════════════════════════════════════════════════════════════════════════════
\newpage
\section{Semantic Hazard Classifier}
% ══════════════════════════════════════════════════════════════════════════════

\subsection{HSV Threshold Rules}

The classifier evaluates a $3{\times}4$ patch grid extracted from the OAK-D
RGB stream. For each patch, hue, saturation, brightness, and edge density
are tested against the following thresholds (all HSV values in OpenCV scale:
$H \in [0,180]$, $S,V \in [0,255]$):

\begin{center}
\begin{tabular}{l c c c c c}
  \toprule
  Class & $H$ range & $S$ threshold & $V$ threshold
        & Edge density $\rho$ & Patch coverage \\
  \midrule
  FIRE    & $<20°$ or $>160°$ & $>120$ & $>120$ & ---   & $>12\%$ \\
  WATER   & $90°$--$130°$      & $>20$  & $>80$  & ---   & $>18\%$ \\
  SMOKE   & any                & $<30$  & $60$--$170$ & --- & $>55\%$ \\
  HAZMAT cue & $18°$--$36°$   & $>100$ & $>100$ & ---   & $>15\%$ \\
  DEBRIS  & any                & any    & any    & $>0.22$ & --- \\
  DARK    & any                & any    & $\bar V < 45$ & --- & mean brightness \\
  \bottomrule
\end{tabular}
\end{center}

\subsection{Edge Density Formula}

For a patch region-of-interest (ROI), edge density $\rho$ is computed
using the Canny operator with thresholds $\tau_1 = 50$, $\tau_2 = 150$:

\begin{equation}
  \rho_{\text{edge}}
  = \frac{\left|\left\{p \in \text{ROI} : |\nabla I(p)| > \tau\right\}\right|}{|\text{ROI}|}
  \tag{18}
\end{equation}

Debris, cables, and structural collapse produce $\rho > 0.22$;
smooth clean floors produce $\rho < 0.10$.

\subsection{HAZMAT Detection: YOLOv3-tiny Decoding and NMS}

DeepHAZMAT uses a YOLOv3-tiny backbone. The predicted bounding box
centre, width, and height are decoded from raw network outputs
$(t_x, t_y, t_w, t_h)$ as:

\begin{equation}
  b_x = \sigma(t_x) + c_x, \quad
  b_y = \sigma(t_y) + c_y, \quad
  b_w = p_w e^{t_w}, \quad
  b_h = p_h e^{t_h}
  \tag{19}
\end{equation}

where $(c_x, c_y)$ is the grid cell offset, $(p_w, p_h)$ are the prior
anchor dimensions, and $\sigma$ is the sigmoid function.

After decoding, Adaptive Non-Maximum Suppression removes redundant boxes:

\begin{equation}
  \text{IoU}(A,B) = \frac{|A \cap B|}{|A \cup B|}
  \tag{20}
\end{equation}

Boxes with $\text{IoU} > \tau_{\text{NMS}} = 0.3$ and lower confidence
are suppressed. Our minimum confidence threshold is $p_o > 0.5$.

\subsection{Semantic Map Encoding}

Hazard labels are encoded as integer cell values in the
\texttt{OccupancyGrid} message:

\begin{center}
\begin{tabular}{c l l}
  \toprule
  Cell value & Hazard class & Behavior \\
  \midrule
  0   & CLEAR                  & Navigate freely \\
  30  & DEBRIS                 & Navigate cautiously \\
  50  & SMOKE                  & Flag for logging \\
  70--75 & WATER / DARK        & Skip as frontier target \\
  80  & NARROW / DEAD\_END     & Skip as frontier target \\
  85--95 & GLASS / HAZMAT / POTHOLE & Skip + E-stop \\
  100 & FIRE / CLIFF           & Immediate E-stop \\
  \bottomrule
\end{tabular}
\end{center}

% ══════════════════════════════════════════════════════════════════════════════
\newpage
\section{Nav2 DWA Local Planner}
% ══════════════════════════════════════════════════════════════════════════════

The Dynamic Window Approach (DWA) samples feasible velocity commands and
selects the one that best balances goal heading, obstacle clearance,
and forward speed.

\subsection{Velocity Constraints}

The dynamic window restricts commands to velocities reachable within
one control cycle $\Delta t$:

\begin{equation}
  v \in [v_{\min}, v_{\max}], \qquad
  \omega \in [-\omega_{\max}, \omega_{\max}]
  \tag{21}
\end{equation}

\begin{equation}
  v \in [v_c - a_v \Delta t,\; v_c + a_v \Delta t], \qquad
  \omega \in [\omega_c - a_\omega \Delta t,\; \omega_c + a_\omega \Delta t]
  \tag{22}
\end{equation}

Safety constraint --- cannot command velocities that cannot brake before
hitting an obstacle at distance $d_{\text{obstacle}}$:

\begin{equation}
  v \leq \sqrt{2 \cdot d_{\text{obstacle}} \cdot a_{v,\text{brake}}}
  \tag{23}
\end{equation}

\subsection{Objective Function}

\begin{equation}
  G(v,\omega) = \sigma\!\left(
    \alpha \cdot \text{heading}(v,\omega)
    + \beta  \cdot \text{clearance}(v,\omega)
    + \gamma \cdot \text{velocity}(v,\omega)
  \right)
  \tag{24}
\end{equation}

where $\sigma$ is a normalisation factor and the three sub-costs are:
\emph{heading} --- alignment with the goal direction;
\emph{clearance} --- distance to nearest obstacle;
\emph{velocity} --- forward speed reward.

\subsection{Arc Trajectory Prediction}

During velocity sampling, DWA predicts robot position using exact
arc integration (same as equation~7):

\begin{equation}
  \begin{bmatrix} x_{k+1} \\ y_{k+1} \\ \theta_{k+1} \end{bmatrix}
  =
  \begin{bmatrix}
    x_k + R\bigl(\sin(\theta_k + \Delta\theta) - \sin\theta_k\bigr) \\
    y_k - R\bigl(\cos(\theta_k + \Delta\theta) - \cos\theta_k\bigr) \\
    \theta_k + \Delta\theta
  \end{bmatrix},
  \qquad
  R = \frac{b}{2} \cdot \frac{\omega_R + \omega_L}{\omega_R - \omega_L}
  \tag{25}
\end{equation}

% ══════════════════════════════════════════════════════════════════════════════
\newpage
\section{3D Differential Drive Kinematics (Uneven Terrain)}
% ══════════════════════════════════════════════════════════════════════════════

For completeness and future M3 hardware testing on uneven surfaces,
the 2D model extends to 6-DOF using ZYX Euler angles.
Control inputs remain the two wheel velocities;
$z$, roll, and pitch are constrained by the terrain.

\subsection{Extended State Vector}

\begin{equation}
  \q(t) = \begin{bmatrix} x\\y\\z\\\phi\\\theta\\\psi \end{bmatrix},
  \qquad
  \bu(t) = \begin{bmatrix}\omega_L\\\omega_R\end{bmatrix}
  \tag{26}
\end{equation}

\begin{center}
\begin{tabular}{>{\(}l<{\)} l}
  \toprule
  \text{Symbol} & Meaning \\
  \midrule
  x,\,y,\,z        & Position in world frame [m] \\
  \phi              & Roll --- rotation about world $x$-axis [rad] \\
  \theta            & Pitch --- rotation about world $y$-axis [rad] \\
  \psi              & Yaw (heading) --- rotation about world $z$-axis [rad] \\
  \omega_L,\,\omega_R & Wheel angular velocities [rad\,s$^{-1}$] \\
  \bottomrule
\end{tabular}
\end{center}

\subsection{Rotation Matrix (ZYX Convention)}

Composing elementary rotations $R_z(\psi)\,R_y(\theta)\,R_x(\phi)$:

\begin{equation}
  R =
  \begin{bmatrix}
    \cps\cth & \cps\sth\sph - \sps\cph & \cps\sth\cph + \sps\sph \\
    \sps\cth & \sps\sth\sph + \cps\cph & \sps\sth\cph - \cps\sph \\
    -\sth    & \cth\sph                & \cth\cph
  \end{bmatrix}
  \tag{27}
\end{equation}

\subsection{World-Frame Translational Velocity}

The robot moves at speed $v$ along its body $x$-axis:

\begin{equation}
  \begin{bmatrix}\dot x \\ \dot y \\ \dot z\end{bmatrix}
  = v(t)\cdot R\,\eone
  = v(t)\begin{bmatrix}\cps\cth \\ \sps\cth \\ -\sth\end{bmatrix}
  \tag{28}
\end{equation}

On flat ground ($\theta=0$) this reduces exactly to the 2D equations.

\subsection{Complete Discrete-Time State Update}

Only $x$, $y$, $\psi$ are numerically integrated;
the remaining three states are read from terrain model $h(x,y)$:

\begin{equation}
  \boxed{
  \begin{aligned}
    x_{k+1}      &= x_k      + v_k\,\cps_k\cth_k\,\Delta t \\
    y_{k+1}      &= y_k      + v_k\,\sps_k\cth_k\,\Delta t \\
    z_{k+1}      &= h(x_{k+1},\,y_{k+1}) \\
    \phi_{k+1}   &= \phi_{\text{surface}}(x_{k+1},\,y_{k+1},\,\psi_{k+1}) \\
    \theta_{k+1} &= \theta_{\text{surface}}(x_{k+1},\,y_{k+1},\,\psi_{k+1}) \\
    \psi_{k+1}   &= \psi_k   + \dot\psi_k\,\Delta t
  \end{aligned}
  }
  \tag{29}
\end{equation}

\noindent\textbf{Singularity warning:}
The Euler angle parameterisation is undefined at $\theta = \pm90°$ (gimbal lock).
For steep terrain, replace with unit quaternion
$\mathbf{q} = (q_0,q_1,q_2,q_3)^\top$ and use
$\dot{\mathbf{q}} = \tfrac{1}{2}\,\Xi(\mathbf{q})\,\bm\omega_b$.

% ══════════════════════════════════════════════════════════════════════════════
\end{document}
